% SEGMENT OLP-0618-S01
% Part: methods
% Chapter: induction
% Section: structural-induction

\documentclass[../../../include/open-logic-section]{subfiles}

\begin{document}

\olfileid{mth}{ind}{sti}

\olsection{கட்டமைப்புத் தொகுத்தறிதல்}

இதுவரை எல்லா இயல் எண்களையும் பற்றிய முடிவுகளை
நிறுவத் தொகுத்தறிதலைப் பயன்படுத்தினோம். ஆனால்
அதற்கொப்பான கொள்கையை, தொகுத்தறிதல் முறையில்
வரையறுக்கப்பட்ட கணத்தின் எல்லா !!{element}s பற்றிய
முடிவுகளுக்கும் நேரடியாகப் பயன்படுத்தலாம். வரையறுக்கப்பட்ட
பொருள்களின் கட்டமைப்பைச் சார்ந்திருப்பதால் இது பல
நேரங்களில் \emph{கட்டமைப்புத்} தொகுத்தறிதல் எனப்படுகிறது.

% SEGMENT OLP-0618-S02
பொதுவாக, தொகுத்தறிதல் வரையறை (a)~கணத்தின்
``தொடக்க'' !!{element}s பட்டியலையும்,
(b)~ஏற்கெனவே உள்ள !!{element}s இலிருந்து
புதிய உறுப்புகளை உருவாக்கும் செயல்களின்
பட்டியலையும் தருகிறது. எடுத்துக்காட்டாக, ஒழுங்கான
சொற்களில் தொடக்கப் பொருள்கள் எழுத்துகள்.
ஒரே ஒரு செயல் உள்ளது; அது:
\begin{align*}
  o(s_1, s_2) = & [s_1 \circ s_2]
\end{align*}
இயல் எண்கள்~$\Nat$ தாமே ஒரு தொகுத்தறிதல்
வரையறையால் தரப்படுவதாகவும் கருதலாம்: தொடக்கப்
பொருள்~$0$; செயல் அடுத்தெண் சார்பு~$x + 1$.

% SEGMENT OLP-0618-S03
தொகுத்தறிதல் முறையில் வரையறுக்கப்பட்ட கணத்தின்
எல்லா உறுப்புகளையும் பற்றி ஏதாவது ஒன்றை,
அதாவது அதன் ஒவ்வோர் !!{element} க்கும்
ஒரு பண்பு~$P$ உண்டு என்று, நிறுவ
வேண்டுமானால்:
\begin{enumerate}
\item தொடக்கப் பொருள்களுக்கு~$P$ உண்டு என்று நிறுவுக;
\item ஒவ்வொரு செயல்~$o$ க்கும், அதன் உள்ளீடுகளுக்கு
  $P$ உண்டு எனில், அதன் முடிவுக்கும் அது உண்டு
  என்று நிறுவுக.
\end{enumerate}
எடுத்துக்காட்டாக, எல்லா ஒழுங்கான சொற்களையும்
பற்றி ஒன்றை நிறுவ, அது எல்லா எழுத்துகளுக்கும்
மெய் என்றும், $s_1$, $s_2$ ஒவ்வொன்றுக்கும்
மெய் எனில் $[s_1 \circ s_2]$ க்கும்
மெய் என்றும் நிறுவுவோம்.

% SEGMENT OLP-0618-S04
\begin{prop}
  எந்த ஒழுங்கான சொல்~$t$ இலும் $[$
  களின் எண்ணிக்கை $]$ களின் எண்ணிக்கைக்குச் சமம்.
\end{prop}

% SEGMENT OLP-0618-S05
\begin{proof}
கட்டமைப்புத் தொகுத்தறிதலைப் பயன்படுத்துவோம்.
ஒழுங்கான சொற்கள் தொகுத்தறிதல் முறையில்
வரையறுக்கப்பட்டுள்ளன: தொடக்கப் பொருள்கள்
எழுத்துகள்; பழைய ஒழுங்கான சொற்களிலிருந்து
புதியவற்றை உருவாக்கும் செயல் $o$.
\begin{enumerate}
\item ஒவ்வோர் எழுத்துக்கும் கூற்று மெய்:
  தனி எழுத்தில் $[$ களின் எண்ணிக்கை~$0$;
  $]$ களின் எண்ணிக்கையும்~$0$.
\item $s_1$ இல் $[$ களின் எண்ணிக்கை
  $]$ களின் எண்ணிக்கைக்குச் சமம் என்றும்,
  $s_2$ க்கும் அதேபோல் என்றும் என்க.
  $o(s_1, s_2)$ இல், அதாவது
  $[s_1 \circ s_2]$ இல் உள்ள $[$ களின்
  எண்ணிக்கை, $s_1$, $s_2$ இரண்டிலும் உள்ள
  $[$ களின் எண்ணிக்கைகளின் கூட்டுத்தொகையுடன்
  ஒன்று சேர்த்தது. $o(s_1, s_2)$ இல் உள்ள
  $]$ களின் எண்ணிக்கையும், $s_1$, $s_2$
  இரண்டிலும் உள்ள $]$ களின் எண்ணிக்கைகளின்
  கூட்டுத்தொகையுடன் ஒன்று சேர்த்தது. ஆகவே
  $o(s_1, s_2)$ இல் $[$ களின் எண்ணிக்கை,
  $o(s_1,s_2)$ இல் $]$ களின் எண்ணிக்கைக்குச்
  சமம்.
\end{enumerate}
\end{proof}

% SEGMENT OLP-0618-S06
\begin{prob}
  எந்த ஒழுங்கான சொல்லும்~$]$ கொண்டு
  தொடங்காது என்பதைக் கட்டமைப்புத்
  தொகுத்தறிதலால் நிறுவுக.
\end{prob}

% SEGMENT OLP-0618-S07
கட்டமைப்புத் தொகுத்தறிதலால் இன்னொரு நிறுவலைத்
தருவோம். குறிகளால் ஆன ஒரு தொடர்~$t$ இன்
முறையான தொடக்கப் பகுதி என்பது,
இடமிருந்து வாசிக்கும்போது குறி குறியாக~$t$
உடன் ஒத்திருந்து, $t$ ஐவிடக் குறைவான
நீளம் உடைய ஏதேனும் ஒரு தொடர்~$s$.
ஆகவே $[a \circ {}$ என்பது $[a \circ b]$
இன் முறையான தொடக்கப் பகுதி. ஆனால்
$[b \circ {}$ அப்படியல்ல (இரண்டாவது
குறியில் வேறுபடுகிறது); $[a \circ b]$
உம் அப்படியல்ல (அதே நீளம் உடையது).

% SEGMENT OLP-0618-S08
\begin{prop}\ollabel{prop:initial}
  ஒழுங்கான சொல்~$t$ இன் ஒவ்வொரு முறையான
  தொடக்கப் பகுதியிலும் $]$ களைவிட $[$
  கள் அதிகமாக இருக்கும்.
\end{prop}

% SEGMENT OLP-0618-S09
\begin{proof}
  ~ $t$ மீது தொகுத்தறிதல் செய்கிறோம்:
  \begin{enumerate}
  \item $t$ ஒரு தனி எழுத்து: அப்போது $t$
    க்கு முறையான தொடக்கப் பகுதிகள் இல்லை.
  \item சில ஒழுங்கான சொற்கள் $s_1$,~$s_2$
    க்கு $t = [s_1 \circ s_2]$ என்க.
    $r$ என்பது $t$ இன் முறையான தொடக்கப்
    பகுதி எனில், பல வாய்ப்புகள் உள்ளன:
    \begin{enumerate}
    \item $r$ என்பது வெறும் $[$: அப்போது
      $r$ இல் $]$ களைவிட ஒரு $[$ அதிகம்.
    \item $r$ என்பது $[r_1$; இங்கே $r_1$
      என்பது~$s_1$ இன் முறையான தொடக்கப்
      பகுதி. $s_1$ ஒழுங்கான சொல் என்பதால்,
      தொகுத்தறிதல் எடுகோளின்படி $r_1$ இல்
      $]$ களைவிட $[$ கள் அதிகம்;
      $[r_1$ இலும் அப்படியே.
    \item $r$ என்பது $[s_1$ அல்லது
      $[s_1 \circ {}$: முந்தைய முடிவின்படி,
      ~ $s_1$ இல் $[$, $]$ களின்
      எண்ணிக்கைகள் சமம். ஆகவே $[s_1$
      அல்லது $[s_1 \circ {}$ இல் $]$
      களைவிட $[$ ஒன்று அதிகம்.
    \item $r$ என்பது $[s_1 \circ r_2$;
      இங்கே $r_2$ என்பது~$s_2$ இன் முறையான
      தொடக்கப் பகுதி. தொகுத்தறிதல் எடுகோளின்படி
      $r_2$ இல் $]$ களைவிட $[$ கள்
      அதிகம். முந்தைய முடிவின்படி~$s_1$
      இல் $[$, $]$ களின் எண்ணிக்கைகள்
      சமம். ஆகவே $[s_1 \circ r_2$ இல்
      $]$ களைவிட $[$ கள் அதிகம்.
    \item $r$ என்பது $[s_1 \circ s_2$:
      முந்தைய முடிவின்படி $s_1$ இல் $[$,
      $]$ களின் எண்ணிக்கைகள் சமம்;~$s_2$
      க்கும் அப்படியே. ஆகவே $[s_1 \circ s_2$
      இல் $]$ களைவிட $[$ ஒன்று அதிகம்.
    \end{enumerate}
  \end{enumerate}
\end{proof}

\end{document}
